\section{Circuit Complexity --- Strict Size}
%%% generated by the blueprint pipeline from TCSlib.Complexity.CircuitComplexity.StrictSize
\subsection*{Circuit Complexity --- Strict Size: Introduction}

Declarations of the Lean module \texttt{TCSlib.Complexity.CircuitComplexity.StrictSize}.

\subsection*{Circuit Complexity --- Strict Size: Results}

\begin{definition}[Three-gate sealing block]
\lean{BoolCircuit.sealTriple}
\label{BoolCircuit.sealTriple}
\leanok
\uses{BoolCircuit.DAGGate}
Given natural numbers $u$, $m$, and $o$, define a list of three gates, in order: a
negation gate reading $u$, a disjunction gate reading $u$ and $m$, and a conjunction
gate reading $o$ and $m+1$.
\end{definition}

\begin{definition}[One-step gate sealing]
\lean{BoolCircuit.sealStep}
\label{BoolCircuit.sealStep}
\leanok
\uses{BoolCircuit.DAGGate, BoolCircuit.sealTriple}
Given a natural number $n$, a state consisting of a gate list and a natural-number
counter, and a natural number $u$, append to the gate list the three-gate sealing block
with arguments $u$, $n$ plus the current list length, and the current counter. Set the
new counter to $n$ plus the current list length plus two.
\end{definition}

\begin{definition}[Iterated gate sealing]
\lean{BoolCircuit.sealGates}
\label{BoolCircuit.sealGates}
\leanok
\uses{BoolCircuit.DAGGate, BoolCircuit.sealStep}
Given a natural number $n$, a list of natural numbers, and a state consisting of a gate
list and a natural-number counter, apply the one-step gate-sealing operation
successively to each number in the list, starting from the given state.
\end{definition}

\begin{theorem}[Fanout in a three-gate construction]
\lean{BoolCircuit.fanoutIn_sealTriple}
\label{BoolCircuit.fanoutIn_sealTriple}
\leanok
\uses{BoolCircuit.fanoutIn, BoolCircuit.sealTriple}
\difficulty{3}
Let $u,N,o,v$ be nonnegative integers with $u<N$ and $o<N$. Form three gates in order: a
NOT gate reading $u$, an OR gate reading $u$ and $N$, and an AND gate reading $o$ and
$N+1$. The fanout of $v$ among these gates is $2$ if $v=u$, plus $1$ if $v=N$, plus $1$
if $v=o$, plus $1$ if $v=N+1$, with each contribution omitted when its stated equality
fails.
\end{theorem}

\begin{theorem}[Strictness of the sealing gates]
\lean{BoolCircuit.sealTriple_strict}
\label{BoolCircuit.sealTriple_strict}
\leanok
\uses{BoolCircuit.DAGGate.Strict, BoolCircuit.sealTriple}
\difficulty{4}
Let $u,N,o$ be natural numbers with $u<N$ and $o<N$. Each of the three gates consisting,
in order, of a NOT gate reading $u$, an OR gate reading $u$ and $N$, and an AND gate
reading $o$ and $N+1$ has distinct arguments; the NOT gate has one argument and each
other gate has two.
\end{theorem}

\begin{theorem}[Acyclicity of the sealing gates]
\lean{BoolCircuit.sealTriple_acyclic}
\label{BoolCircuit.sealTriple_acyclic}
\leanok
\uses{BoolCircuit.GatesAcyclic, BoolCircuit.gatesAcyclic_cons, BoolCircuit.gatesAcyclic_nil, BoolCircuit.sealTriple}
\difficulty{3}
Let $u,N,o$ be natural numbers with $u<N$ and $o<N$. The three gates consisting, in
order, of a NOT gate reading $u$, an OR gate reading $u$ and $N$, and an AND gate
reading $o$ and $N+1$ form an acyclic gate list relative to $N$: gate $i$ reads only
vertices with indices less than $N+i$.
\end{theorem}

\begin{theorem}[Preservation of an existing circuit value]
\lean{BoolCircuit.vertexValue_sealStep}
\label{BoolCircuit.vertexValue_sealStep}
\leanok
\uses{BoolCircuit.DAGGate, BoolCircuit.DAGGate.eval, BoolCircuit.DAGGate.eval_pair, BoolCircuit.sealStep, BoolCircuit.sealTriple, BoolCircuit.vertexValue, BoolCircuit.vertexValue_append, BoolCircuit.vertexValue_last}
\difficulty{5}
Let $n$ be a nonnegative integer, let $G$ be a finite list of Boolean gates, let $x$
assign a Boolean value to each of the first $n$ inputs, and let $u,o$ be gate-value
indices satisfying $u<n+|G|$ and $o<n+|G|$. Appending the three gates that negate the
value at $u$, take its disjunction with the value at index $n+|G|$, and then conjoin the
value at $o$ with the value at index $n+|G|+1$ leaves the value at output index
$n+|G|+2$ equal to the original value at index $o$.
\end{theorem}

\begin{definition}[Specification for sealed gates]
\lean{BoolCircuit.SealSpec}
\label{BoolCircuit.SealSpec}
\leanok
\uses{BoolCircuit.DAGGate, BoolCircuit.DAGGate.Strict, BoolCircuit.GatesAcyclic, BoolCircuit.vertexValue}
For a natural number $n$, an original gate list $G$, an output vertex $o$, and a
resulting pair $(H,r)$ of a gate list and a natural number, the specification requires
that $H$ extend $G$ by a list of gates each having distinct arguments and the prescribed
arity: one for negation gates and two for other gates. The resulting gate list must be
acyclic, $r$ must be a vertex of that list, and for every Boolean assignment to the $n$
input vertices, the value at $r$ in the resulting list must equal the value at $o$ in
the original list.
\end{definition}

\begin{theorem}[Length after sealing gates]
\lean{BoolCircuit.sealGates_length}
\label{BoolCircuit.sealGates_length}
\leanok
\uses{BoolCircuit.DAGGate, BoolCircuit.sealGates, BoolCircuit.sealStep, BoolCircuit.sealTriple}
\difficulty{4}
Let $n\in\mathbb{N}$, let $L$ be a list of indices, and let $s$ be a pair consisting of
a gate list and a natural number. The gate list produced by sealing the indices in $L$
starting from $s$ has length equal to the original gate-list length plus three times the
length of $L$.
\end{theorem}

\begin{theorem}[Specification of gate sealing]
\lean{BoolCircuit.sealGates_spec}
\label{BoolCircuit.sealGates_spec}
\leanok
\uses{BoolCircuit.DAGGate, BoolCircuit.GatesAcyclic, BoolCircuit.SealSpec, BoolCircuit.gatesAcyclic_append, BoolCircuit.sealGates, BoolCircuit.sealStep, BoolCircuit.sealTriple, BoolCircuit.sealTriple_acyclic, BoolCircuit.sealTriple_strict, BoolCircuit.vertexValue_sealStep}
\difficulty{6}
Let $n$ be a natural number, let $L$ be a list of vertex indices, let $G$ be a list of
gates, and let $o$ be a vertex index. Suppose $G$ is acyclic, $o<n+|G|$, and every index
in $L$ is less than $n+|G|$. Sealing the indices in $L$ starting from $(G,o)$ produces a
gate list and output index such that the new list extends $G$ by gates with distinct
arguments and the prescribed argument count (one for a NOT gate and two otherwise), the
new list is acyclic, the new output is a valid vertex, and its value agrees with the
value of vertex $o$ in $G$ on every Boolean input in $\{0,1\}^n$.
\end{theorem}

\begin{theorem}[Readability after gate sealing]
\lean{BoolCircuit.sealGates_read}
\label{BoolCircuit.sealGates_read}
\leanok
\uses{BoolCircuit.DAGGate, BoolCircuit.fanoutIn, BoolCircuit.fanoutIn_append, BoolCircuit.fanoutIn_sealTriple, BoolCircuit.sealGates, BoolCircuit.sealTriple}
\difficulty{5}
Let $n$ be a natural number, let $L$ be a list of vertex indices, let $G$ be a list of
gates, and let $o$ be a vertex index. Suppose $o<n+|G|$, every index in $L$ is less than
$n+|G|$, and every vertex $v<n+|G|$ with $v\ne o$ and $v\notin L$ is read by at least
one gate in $G$. After sealing the indices in $L$ starting from $(G,o)$, every vertex
below $n$ plus the resulting gate count, except the resulting output vertex, is read by
at least one resulting gate.
\end{theorem}

\begin{theorem}[Fanout bound after sealing gates]
\lean{BoolCircuit.sealGates_fanout}
\label{BoolCircuit.sealGates_fanout}
\leanok
\uses{BoolCircuit.DAGGate, BoolCircuit.GatesAcyclic, BoolCircuit.fanoutIn, BoolCircuit.fanoutIn_append, BoolCircuit.fanoutIn_eq_zero_of_le, BoolCircuit.fanoutIn_sealTriple, BoolCircuit.gatesAcyclic_append, BoolCircuit.sealGates, BoolCircuit.sealTriple, BoolCircuit.sealTriple_acyclic}
\difficulty{6}
Fix $n\in\mathbb{N}$. Let $gs$ be an acyclic gate list, meaning that every argument
index at position $i$ is less than $n+i$, and let $o<n+|gs|$. Let $L$ be a list of
distinct indices, each less than $n+|gs|$, with $o\notin L$. Suppose every index in $L$
has fanout zero in $gs$, the fanout of $o$ in $gs$ is at most one, and every index has
fanout at most two in $gs$. Then after sealing the indices in $L$ starting from
$(gs,o)$, every index has fanout at most two in the resulting gate list.
\end{theorem}

\begin{definition}[Dead vertices of a circuit]
\lean{BoolCircuit.DAGCircuit.deadVertices}
\label{BoolCircuit.DAGCircuit.deadVertices}
\leanok
\uses{BoolCircuit.DAGCircuit, BoolCircuit.DAGCircuit.fanout, BoolCircuit.DAGCircuit.size}
Let $C$ be a Boolean circuit with $n$ input vertices, gates in topological order, and a
designated output vertex. Its dead-vertex list consists of the vertices $v$ with $0 \le
v < n + \mathrm{length}(C.\mathrm{gates})$ that are not the output and are read by no
gate.
\end{definition}

\begin{theorem}[Characterization of dead vertices]
\lean{BoolCircuit.DAGCircuit.mem_deadVertices}
\label{BoolCircuit.DAGCircuit.mem_deadVertices}
\leanok
\uses{BoolCircuit.DAGCircuit, BoolCircuit.DAGCircuit.deadVertices, BoolCircuit.DAGCircuit.fanout, BoolCircuit.DAGCircuit.size}
\difficulty{1}
Let $C$ be a directed acyclic circuit with $n$ input vertices, gates in topological
order, and a valid output vertex. A vertex belongs to the list of dead vertices exactly
when it is less than the circuit size, is not the output vertex, and is read by no gate.
\end{theorem}

\begin{definition}[Sealing a circuit]
\lean{BoolCircuit.DAGCircuit.seal}
\label{BoolCircuit.DAGCircuit.seal}
\leanok
\uses{BoolCircuit.DAGCircuit, BoolCircuit.DAGCircuit.deadVertices, BoolCircuit.DAGCircuit.mem_deadVertices, BoolCircuit.sealGates, BoolCircuit.sealGates_spec}
Given a circuit with $n$ input vertices, process its dead vertices in order. For each
dead vertex $u$, append three gates: a NOT gate reading $u$, an OR gate reading $u$ and
the next new vertex, and an AND gate reading the current output and the vertex produced
by the OR gate; update the output to the vertex produced by this AND gate. The resulting
circuit is the sealed circuit.
\end{definition}

\begin{theorem}[Specification of circuit sealing]
\lean{BoolCircuit.seal_spec}
\label{BoolCircuit.seal_spec}
\leanok
\uses{BoolCircuit.DAGCircuit.deadVertices, BoolCircuit.DAGCircuit.mem_deadVertices, BoolCircuit.SealSpec, BoolCircuit.sealGates, BoolCircuit.sealGates_spec}
\difficulty{1}
Let $C$ be a finite acyclic Boolean circuit with input indices $0,\ldots,n-1$, a valid
output vertex, and gates whose arguments are distinct and have the prescribed count: one
for a NOT gate and two otherwise. Seal every vertex index below the circuit size that is
not the output and has fanout zero, starting from the gates and output of $C$. The
resulting gates extend the original list by gates with distinct arguments and the
prescribed argument counts, remain acyclic, and have a valid output whose value agrees
with the original circuit's output on every Boolean input.
\end{theorem}

\begin{theorem}[Evaluation preserved by sealing]
\lean{BoolCircuit.DAGCircuit.seal_eval}
\label{BoolCircuit.DAGCircuit.seal_eval}
\leanok
\uses{BoolCircuit.DAGCircuit.eval, BoolCircuit.DAGCircuit.seal, BoolCircuit.seal_spec}
\difficulty{1}
Let $C$ be a finite Boolean circuit with $n$ input vertices, gates in topological order,
and a designated output vertex. For every Boolean input assignment $x$, the sealed
circuit has the same output value as $C$.
\end{theorem}

\begin{theorem}[Size bound for the sealed circuit]
\lean{BoolCircuit.DAGCircuit.seal_size_le}
\label{BoolCircuit.DAGCircuit.seal_size_le}
\leanok
\uses{BoolCircuit.DAGCircuit.deadVertices, BoolCircuit.DAGCircuit.fanout, BoolCircuit.DAGCircuit.seal, BoolCircuit.DAGCircuit.size, BoolCircuit.sealGates, BoolCircuit.sealGates_length}
\difficulty{4}
Let $C$ be a finite Boolean circuit with $n$ input vertices, gates in topological order,
and a designated output vertex. The size of its sealed circuit is at most four times the
size of $C$.
\end{theorem}

\begin{theorem}[Strictness after sealing]
\lean{BoolCircuit.DAGCircuit.seal_isStrict}
\label{BoolCircuit.DAGCircuit.seal_isStrict}
\leanok
\uses{BoolCircuit.DAGCircuit.IsStrict, BoolCircuit.DAGCircuit.deadVertices, BoolCircuit.DAGCircuit.fanout, BoolCircuit.DAGCircuit.mem_deadVertices, BoolCircuit.DAGCircuit.seal, BoolCircuit.DAGGate.Strict, BoolCircuit.sealGates, BoolCircuit.sealGates_read, BoolCircuit.seal_spec}
\difficulty{5}
Let $C$ be a finite Boolean circuit with $n$ input vertices, gates in topological order,
and a designated output vertex. If every gate of $C$ has distinct arguments and has one
argument when it is a NOT gate and two otherwise, then the sealed circuit has those same
gate properties, and every vertex other than its output is read by at least one gate.
\end{theorem}

\begin{theorem}[Fanout bound after sealing]
\lean{BoolCircuit.DAGCircuit.seal_hasFanoutTwo}
\label{BoolCircuit.DAGCircuit.seal_hasFanoutTwo}
\leanok
\uses{BoolCircuit.DAGCircuit.HasFanoutTwo, BoolCircuit.DAGCircuit.deadVertices, BoolCircuit.DAGCircuit.fanout, BoolCircuit.DAGCircuit.mem_deadVertices, BoolCircuit.DAGCircuit.seal, BoolCircuit.sealGates_fanout}
\difficulty{1}
Let $C$ be a finite Boolean circuit with $n$ input vertices, gates in topological order,
and a designated output vertex. If every vertex of $C$ is read by at most two gates and
its output vertex is read by at most one gate, then every vertex of the sealed circuit
is read by at most two gates.
\end{theorem}

\begin{definition}[Strictification of a circuit]
\lean{BoolCircuit.DAGCircuit.strictify}
\label{BoolCircuit.DAGCircuit.strictify}
\leanok
\uses{BoolCircuit.DAGCircuit, BoolCircuit.DAGCircuit.deconst, BoolCircuit.DAGCircuit.seal}
Let $C$ be a circuit with $n>0$ input vertices. First prepend three gates, a NOT gate
reading vertex $0$, an OR gate reading vertices $0$ and $n$, and an AND gate reading
vertices $0$ and $n$; shift every original vertex index at least $n$ upward by $3$,
applying the corresponding gate transformation to each original gate and shifting the
output index as well. Then seal the resulting circuit by processing its dead vertices.
\end{definition}

\begin{theorem}[Evaluation preserved by strictification]
\lean{BoolCircuit.DAGCircuit.strictify_eval}
\label{BoolCircuit.DAGCircuit.strictify_eval}
\leanok
\uses{BoolCircuit.DAGCircuit.IsFaninTwo, BoolCircuit.DAGCircuit.deconst_eval, BoolCircuit.DAGCircuit.eval, BoolCircuit.DAGCircuit.seal_eval, BoolCircuit.DAGCircuit.strictify}
\difficulty{3}
Let $C$ be a finite Boolean circuit with $n>0$ input vertices, gates in topological
order, and a designated output vertex. Suppose each gate has distinct arguments, every
NOT gate has one argument, and every gate has at most two arguments. For every Boolean
input assignment $x$, the strictified circuit has the same output value as $C$.
\end{theorem}

\begin{theorem}[Strictness after strictification]
\lean{BoolCircuit.DAGCircuit.strictify_isStrict}
\label{BoolCircuit.DAGCircuit.strictify_isStrict}
\leanok
\uses{BoolCircuit.DAGCircuit.IsFaninTwo, BoolCircuit.DAGCircuit.IsStrict, BoolCircuit.DAGCircuit.deconst, BoolCircuit.DAGCircuit.deconst_strict, BoolCircuit.DAGCircuit.seal_isStrict, BoolCircuit.DAGCircuit.strictify}
\difficulty{1}
Let $C$ be a finite Boolean circuit with $n>0$ input vertices, gates in topological
order, and a designated output vertex. Suppose each gate has distinct arguments, every
NOT gate has one argument, and every gate has at most two arguments. Then the
strictified circuit has distinct arguments at every gate, each NOT gate has exactly one
argument and each other gate exactly two arguments, and every vertex other than its
output is read by at least one gate.
\end{theorem}

\begin{theorem}[Size bound for strictification]
\lean{BoolCircuit.DAGCircuit.strictify_size_le}
\label{BoolCircuit.DAGCircuit.strictify_size_le}
\leanok
\uses{BoolCircuit.DAGCircuit.deconst, BoolCircuit.DAGCircuit.deconst_size, BoolCircuit.DAGCircuit.seal_size_le, BoolCircuit.DAGCircuit.size, BoolCircuit.DAGCircuit.strictify}
\difficulty{3}
Let $C$ be a directed acyclic Boolean circuit with $n>0$ input vertices, and define its
size to be $n$ plus its number of gates. The strictification of $C$ has size at most
$4\,\mathrm{size}(C)+12$.
\end{theorem}

\begin{theorem}[Strict circuit conversion]
\lean{BoolCircuit.DAGCircuit.exists_isStrict}
\label{BoolCircuit.DAGCircuit.exists_isStrict}
\leanok
\uses{BoolCircuit.DAGCircuit, BoolCircuit.DAGCircuit.IsFaninTwo, BoolCircuit.DAGCircuit.IsStrict, BoolCircuit.DAGCircuit.eval, BoolCircuit.DAGCircuit.size, BoolCircuit.DAGCircuit.strictify, BoolCircuit.DAGCircuit.strictify_eval, BoolCircuit.DAGCircuit.strictify_isStrict, BoolCircuit.DAGCircuit.strictify_size_le}
\difficulty{1}
Let $n$ be a positive integer, and let $C$ be a directed acyclic Boolean circuit on $n$
inputs in which every gate has distinct arguments and at most two arguments, and every
NOT gate has exactly one argument. There is a circuit $C'$ on the same inputs that
computes the same Boolean function, in which every gate has distinct arguments and
exactly one argument for a NOT gate or two for an AND or OR gate, and every vertex other
than the output has positive fanout; moreover, $C'$ has size at most $4$ times the size
of $C$ plus $12$.
\end{theorem}

\begin{theorem}[Equivalent circuit representation conditions]
\lean{BoolCircuit.DAGCircuit.exists_isStrict_iff}
\label{BoolCircuit.DAGCircuit.exists_isStrict_iff}
\leanok
\uses{BoolCircuit.DAGCircuit, BoolCircuit.DAGCircuit.IsFaninTwo, BoolCircuit.DAGCircuit.IsStrict, BoolCircuit.DAGCircuit.IsStrict.isFaninTwo, BoolCircuit.DAGCircuit.eval, BoolCircuit.DAGCircuit.exists_isStrict}
\difficulty{4}
Let $n$ be a positive integer and let $f$ be a Boolean function on $n$ inputs. There
exists a directed acyclic Boolean circuit on those inputs, with distinct arguments at
each gate, at most two arguments per gate, and exactly one argument at every NOT gate,
computing $f$, if and only if there exists such a circuit computing $f$ in which every
gate has exactly one argument for a NOT gate or two for an AND or OR gate and every
vertex other than the output has positive fanout.
\end{theorem}

\begin{definition}[Strictification of a circuit family]
\lean{BoolCircuit.DAGCircuitFamily.strictify}
\label{BoolCircuit.DAGCircuitFamily.strictify}
\leanok
\uses{BoolCircuit.DAGCircuit.strictify, BoolCircuit.DAGCircuitFamily}
Let $C$ be a family assigning a circuit with $n$ input vertices to each natural number
$n$. Define a new family that uses $C$'s circuit unchanged at $n=0$ and uses the
strictification of $C$'s circuit at every positive $n$.
\end{definition}

\begin{theorem}[Evaluation preserved by strictification]
\lean{BoolCircuit.DAGCircuitFamily.strictify_eval}
\label{BoolCircuit.DAGCircuitFamily.strictify_eval}
\leanok
\uses{BoolCircuit.DAGCircuit.eval, BoolCircuit.DAGCircuit.strictify_eval, BoolCircuit.DAGCircuitFamily, BoolCircuit.DAGCircuitFamily.HasFaninTwo, BoolCircuit.DAGCircuitFamily.strictify}
\difficulty{4}
Let $C$ be a family of Boolean circuits such that every circuit in the family is well
formed and every gate has at most two arguments. For every input length $n$ and every
Boolean input of length $n$, the circuit produced by strictifying $C_n$ has the same
output as $C_n$.
\end{theorem}

\begin{definition}[Languages computed with strict size bound]
\lean{Language.InStrictSIZE}
\label{Language.InStrictSIZE}
\leanok
\uses{BoolCircuit.DAGCircuit.IsFaninTwo, BoolCircuit.DAGCircuit.IsStrict, BoolCircuit.DAGCircuit.size, BoolCircuit.DAGCircuitFamily, BoolCircuit.DAGCircuitFamily.language}
A language $L$ over the Boolean alphabet belongs to the strict size class for a function
$T:\mathbb{N}\to\mathbb{N}$ if there is a family of Boolean circuits such that the
circuit at input length $0$ has fan-in at most two, every circuit at positive input
length is strict, each circuit at input length $n$ has size at most $T(n)$, and the
family computes exactly $L$.
\end{definition}

\begin{definition}[Strict polynomial-size circuit class]
\lean{Language.InStrictPPoly}
\label{Language.InStrictPPoly}
\leanok
\uses{Language.InStrictSIZE}
A Boolean language belongs to the strict polynomial-size circuit class if there are
natural numbers $a$ and $k$ such that it belongs to the strict size class with bound
$n\mapsto a(n+1)^k$.
\end{definition}

\begin{theorem}[Strict size classes lie in size classes]
\lean{Language.InStrictSIZE.inSIZE}
\label{Language.InStrictSIZE.inSIZE}
\leanok
\uses{BoolCircuit.DAGCircuit.IsStrict.isFaninTwo, Language.InSIZE, Language.InStrictSIZE}
\difficulty{4}
For a function $T : \mathbb{N} \to \mathbb{N}$ and a language $L$ over Boolean strings,
if $L$ is recognized by a family of circuits whose zero-input-length circuit is well
formed with fan-in at most two, whose circuits for every positive input length have
distinct inputs to each gate with exactly one input for NOT gates and exactly two for
AND and OR gates, and in which every vertex other than the output is used by some gate,
with size at most $T(n)$ at every input length $n$, then $L$ is recognized by a family
of well-formed fan-in-two circuits of size at most $T(n)$ at every input length $n$.
\end{theorem}

\begin{theorem}[Monotonicity of strict circuit size bounds]
\lean{Language.InStrictSIZE.mono}
\label{Language.InStrictSIZE.mono}
\leanok
\uses{Language.InStrictSIZE}
\difficulty{1}
For functions $T,T' : \mathbb{N} \to \mathbb{N}$ and a language $L$ over Boolean
strings, if $L$ is recognized by a strict circuit family of size at most $T(n)$ at every
input length $n$, and $T(n) \leq T'(n)$ for every $n$, then $L$ is recognized by a
strict circuit family of size at most $T'(n)$ at every input length $n$.
\end{theorem}

\begin{theorem}[Conversion to strict circuit size]
\lean{Language.InSIZE.inStrictSIZE}
\label{Language.InSIZE.inStrictSIZE}
\leanok
\uses{BoolCircuit.DAGCircuit.strictify_isStrict, BoolCircuit.DAGCircuit.strictify_size_le, BoolCircuit.DAGCircuitFamily.mem_language_iff, BoolCircuit.DAGCircuitFamily.strictify, BoolCircuit.DAGCircuitFamily.strictify_eval, Language.InSIZE, Language.InStrictSIZE}
\difficulty{5}
Let $T : \mathbb{N} \to \mathbb{N}$ and let $L$ be a language of Boolean strings. If $L$
is recognized by a circuit family with fan-in at most two and size at most $T(n)$ on
inputs of length $n$, then it is recognized by a family whose circuit at length zero has
fan-in at most two and whose circuits at every positive input length are strict: each
gate has distinct inputs and has exactly one input if it is a NOT gate and exactly two
otherwise, and every vertex other than the output has positive fanout. The size at input
length $n$ is at most $4T(n)+12$.
\end{theorem}

\begin{theorem}[Strict and unrestricted polynomial size]
\lean{Language.inPPoly_iff_inStrictPPoly}
\label{Language.inPPoly_iff_inStrictPPoly}
\leanok
\uses{Language.InPPoly, Language.InSIZE.inStrictSIZE, Language.InStrictPPoly, Language.InStrictSIZE.inSIZE, Language.InStrictSIZE.mono}
\difficulty{4}
For every Boolean language $L$, recognition by polynomial-size families of directed
acyclic Boolean circuits with fan-in at most two is equivalent to the existence of
natural numbers $a,k$ and a circuit family whose size is at most $a(n+1)^k$ for every
input length $n$, whose circuit at length zero has fan-in at most two, and whose
circuits at every positive input length have gates with distinct inputs and exactly one
input for NOT gates and exactly two for AND and OR gates, with every vertex other than
the output feeding at least one gate.
\end{theorem}
